\addcontentsline{toc}{chapter}{Príloha}
\chapter*{Príloha}
\vspace{-1cm}

\subsection*{Štruktúra priečinka \texttt{compiler}}

\noindent
\begin{minipage}[t]{0.48\textwidth}
\dirtree{%
.1 compiler/.
.2 generate\_tests.sh.
.2 pom.xml.
.2 README.md.
.2 src/.
.3 main/.
.4 antlr4/.
.5 datalog.g4.
.4 java/.
.5 common/.
.6 Constant.java.
.6 FunctionTerm.java.
.6 Parameter.java.
.6 Term.java.
.6 Variable.java.
.5 datalog/.
.6 Atom.java.
.6 Clause.java.
.6 DatalogProgram.java.
.6 Definition.java.
.6 Predicate.java.
.6 Query.java.
.5 ra/.
.6 RAAntiJoin.java.
.6 RADef.java.
.6 RAExpr.java.
.6 RAJoin.java.
.6 RARec.java.
.6 RAUnion.java.
.6 RelationRef.java.
.5 DatalogASTBuilder.java.
.5 DatalogCompiler.java.
.5 MyErrorListener.java.
.5 TranslateDatalogFile.java.
.2 tests/.
.3 000\_turing\_machine.dl.
.3 000\_turing\_machine.ra.
.3 \ldots.
}
\end{minipage}\hfill
\begin{minipage}[t]{0.48\textwidth}
Súbory s príponou \texttt{.dl} sú vstupné testovacie Datalogové programy. Súbory s príponou
\texttt{.ra} obsahujú očakávaný výstup prekladu v relačnej algebre. Testy označené ako
\textit{chybový vstup} (prípona \texttt{\_e}) neobsahujú súbor \texttt{.ra}, pretože ich očakávaným
výsledkom je chybová hláška kompilátora.

\vline\vspace{-0.3cm}

Súbor \texttt{TranslateDatalogFile.java} obsahuje metódu \texttt{main}, ktorá preloží Datalogový
program zo štandardného vstupu a vypíše výsledok na štandardný výstup.

\vline\vspace{-0.3cm}

Súbor \texttt{generate\_tests.sh} je skript, ktorý pre každý Datalogový program v priečinku
\texttt{tests} spustí kompilátor, presmeruje ho na štandardný vstup a výstup presmeruje do
príslušného súboru \texttt{.ra}.

\vline\vspace{-0.3cm}

Súbor \texttt{datalog.g4} obsahuje gramatiku Datalogu, ktorú používa nástroj ANTLR na vygenerovanie
Java kódu pre lexikálnu a syntaktickú analýzu Datalogových programov.

\vline\vspace{-0.3cm}

Súbor \texttt{README.md} obsahuje stručný popis priečinka a návod na spustenie kompilátora a testov.
\end{minipage}

